\documentclass[10pt]{article}
\usepackage[utf8]{inputenc}
\usepackage[T1]{fontenc}
\usepackage{fixltx2e}
\usepackage{amsmath}
\usepackage{amsfonts}
\usepackage{amssymb}
\usepackage[version=4]{mhchem}
\usepackage{stmaryrd}
\usepackage{hyperref}
\hypersetup{colorlinks=true, linkcolor=blue, filecolor=magenta, urlcolor=cyan,}
\urlstyle{same}

\DeclareUnicodeCharacter{25A1}{\ifmmode\square\else{$\square$}\fi}
\DeclareUnicodeCharacter{25EF}{\ifmmode\bigcirc\else{$\bigcirc$}\fi}

\begin{document}
\section*{mathoverflow}
\section*{How to calculate [10\textsuperscript{10}10\textsuperscript{10}10\textsuperscript{-10}10]?}
Asked 14 years, 6 months ago Modified 1 year, 11 months ago Viewed 7k times

\begin{itemize}
  \item How to find an integer part of $10^{10^{10^{10^{10^{-10^{10}}}}}}$ ? It looks like it is slightly above $10^{10^{10}}$.
\end{itemize}

45\\
nt.number-theory\\
computational-number-theory

Share Cite Edit Close Delete Flag\\
edited Oct 27, 2021 at 16:19\\
asked Oct 27, 2011 at 0:18\\
GH from MO\\
Vladimir Reshetnikov

2 Only slightly related, but for amusement, check out \href{http://math.stackexchange.com/questions/72646/}{math.stackexchange.com/questions/72646/}...

\begin{itemize}
  \item Joel David Hamkins Oct 27, 2011 at 0:27
\end{itemize}

4 @quid: It is $-10^{\wedge} 10$, not $(-10)^{\wedge} 10$. - GH from MO Oct 27, 2011 at 0:28

1 @GH: Thanks! So, it was me being confused. Sorry for the noise. - user9072 Oct 27, 2011 at 0:35

3 Just for curiosisty: Was there any context for the appearance of this number? - efs Oct 27, 2021 at 16:26

2 @EFinat-S I experimented with various power towers, just for fun, and then I thought: is it possible to construct a non-trivial problem about power towers that results in a "small" number (i.e. a number whose decimal expansion is practically possible to fully write down), and I came up with this question.

\begin{itemize}
  \item Vladimir Reshetnikov Oct 27, 2021 at 23:05 Delete
\end{itemize}

1 Answer\\
Sorted by: □ Highest score (default)\\
◯ I think the number in question is $10^{10^{10}}+10^{11} \ln ^{4}(10)$ plus a tiny positive number. That is, it starts with a digit 1 , followed by $10^{10}-13$ zeros, then by the string 2811012357389 , then a\\
94 decimal point, and then some garbage (which starts like $4407116278 \ldots$ ).

To see this let $x:=10^{-10^{10}}$, a tiny positive number, and put $c:=\ln (10)$, an important constant. We have

$$
\begin{gathered}
10^{x}=1+c x+O\left(x^{2}\right) \\
10^{10^{x}}=10^{1+c x+O\left(x^{2}\right)}=10+10 c^{2} x+O\left(x^{2}\right) \\
10^{10^{10^{x}}}=10^{10+10 c^{2} x+O\left(x^{2}\right)}=10^{10}+10^{11} c^{3} x+O\left(x^{2}\right) \\
10^{10^{10^{10^{x}}}}=10^{10^{10}+10^{11} c^{3} x+O\left(x^{2}\right)}=10^{10^{10}}+10^{10^{10}} 10^{11} c^{4} x+O\left(x^{2}\right),
\end{gathered}
$$

where $O\left(x^{2}\right)$ means something tiny all the way.\\
In the last expression we have $10^{10^{10}} 10^{11} c^{4} x=10^{11} c^{4}$, which justifies my claim.

35 I'm not convinced the question belongs on MO , but I like the answer too much to vote the question down. - Gerry Myerson Oct 27, 2011 at 4:45

3 @Gerry: I thought the same! :-) - GH from MO Oct 27, 2011 at 14:39

1 Robert Munafo has a "hypercalc", unfortunately not for the windows-environment, so I cannot use it (Actually it is a perl-script). \href{http://mrob.com/pub/perl/index.html}{mrob.com/pub/perl/index.html} cite: "Hypercalc: An unusual calculator program. It represents numbers in a special way allowing the calculation of quantities much larger than tools such as bc, dc, MACSYMA/maxima, Mathematica and Maple, all of which use a bignum library. For example, you can use Hypercalc to determine whether $128^{\wedge} 48^{\wedge} 1024$ is larger than $8^{\wedge} 88^{\wedge} 888$. (...)" (from R. Munafo's site). Someone who could try it with this program? - Gottfried Helms Oct 27, 2011 at 17:41

4 As far as I can tell, hypercalc just uses standard floating point representations of iterated logs. This suggests that it cannot handle the sort of precision necessary here. - S. Carnahan - Oct 28, 2011 at 3:55


\end{document}